\documentclass[12pt,a4paper]{iopart}

% Document information
\newcommand{\DRN}{CAS-PUB-GEN-000000}
\newcommand{\thisversion}{1.0}

% Author information
\input{common/authors/twhyntie.tex}

% Packages and settings for CERN@school document formatting
% Packages and settings for CERN@school document formatting.

% Graphics
\usepackage[xetex]{graphicx}
\usepackage{subcaption}
\usepackage{enumerate}

% Tables
\usepackage{dcolumn}
\newcolumntype{.}{D{.}{.}{2}}

% Layout
\usepackage{layouts}
\usepackage{lscape}
\usepackage{parskip}
%\usepackage[parfill]{parskip}

% Colours
\usepackage{xcolor}
\definecolor{black}{RGB}{0,0,0}
\definecolor{cernblue}{RGB}{0,85,160}

% Mathematics
\usepackage{latexsym}

% Fonts
\usepackage{fontspec,xunicode}
\defaultfontfeatures{Mapping=tex-text,Scale=MatchLowercase}
\setmainfont{Latin Modern Sans}

% Headers and footers
%
% Using the fancyhdr package for the, um, fancy headings...
\usepackage{fancyhdr}
\setlength{\headheight}{16pt}
\setlength{\voffset}{-0.3in}
\pagestyle{fancy}
%
% Redefine the IOP article header settings for the first page.
\makeatletter
\def\ps@myheadings{%
%    \let\@oddfoot\@empty\let\@evenfoot\@empty
    \let\@oddhead\@empty\let\@evenhead\@empty
    \let\@mkboth\@gobbletwo
    \let\sectionmark\@gobble
    \let\subsectionmark\@gobble}
\makeatother
\lhead{\sffamily\rightmark}
\rhead{\sffamily\leftmark}
\lfoot{\sffamily{\DRN-v\thisversion}}
%\lfoot{\sffamily{CERN@school \DRN-v\thisversion}}
\cfoot{}
\rfoot{\sffamily{\thepage}}

% Exercises/questions
\usepackage[listings]{tcolorbox}
\tcbset{before={\par\medskip\pagebreak[0]\noindent},after={\par\medskip}}%

% Coding
\usepackage{listings}
% Listing options
\definecolor{listbacking}{rgb}{0.98,0.98,0.98}
\definecolor{commentgrey}{rgb}{0.60,0.60,0.60}
\definecolor{keywordblue}{rgb}{0.00,0.00,0.50}
\definecolor{stringgreen}{rgb}{0.00,0.50,0.00}
\definecolor{deepred}{rgb}{0.60,0.00,0.00}
\lstset{
  basicstyle=\footnotesize\ttfamily,
  gobble=4,
  float=hbtp,
  emph={__init__},
  emphstyle=\color{deepred},
  showspaces=false,
  showstringspaces=false,
  backgroundcolor=\color{listbacking},
  commentstyle=\color{commentgrey},
  stringstyle=\color{stringgreen},
  keywordstyle=\color{keywordblue},
  numbers=left,
  numbersep=12pt,
  %framexleftmargin=20pt,
  %framexrightmargin=20pt,
  stepnumber=1,
  showlines=true
}

% Acronyms
\usepackage{acronym}

\usepackage{perpage}
\MakePerPage{footnote}

% Background image.
\usepackage{eso-pic}
\newcommand\AtPageLowerRight[1]{\AtPageLowerLeft{%
   \makebox[\paperwidth][r]{#1}}}
\usepackage{ifthen}

% Links
\usepackage[xetex,pdfpagelabels]{hyperref}%
\hypersetup{%
pdfauthor={\theauthorfull},%
colorlinks=true,%
urlcolor=cernblue,%
citecolor=cernblue,%
%linkcolor=cernblue%
linkcolor=black%
}


% CERN@school macros and definitions
\input{common/tools/defs.tex}

\begin{document}

%
% Title
%
% The CERN@school logo
% To insert the CERN@school logo as a header.
\vspace*{-0.8in}
\hspace*{-0.8in}
\includegraphics[height=1.0in]{common/assets/logos/L000001/doc-header.pdf}



\title{%
The CERN@school Programme: Document Index
}
% 
% Author information
\author{\theauthorinit$^{1, \, 2}$}
%
\address{$^1$\theauthoraddressA}
\address{$^2$\theauthoraddressB}
\ead{\mailto{\theauthoremail}}

%-----------------------------------------------------------------------------
% Abstract
\begin{abstract}
CERN@school is a student-led research programme that brings
CERN into the classroom.
This document provides an index of the
documents,
presentations,
events,
code repositories,
and
publications
related to research and other activities associated with the
programme from June 2012 to December 2016.
It therefore effectively serves as a chronicle of the programme
for this period, covering the national expansion of the CERN@school Timepix detector
network, the launch of the Langton Ultimate Cosmic ray Intensity Detector (LUCID),
the TimPix project (part of ESA astronaut's Principia mission),
and the integration of CERN@school with the Monopole and Exotics Detector
at the LHC (MoEDAL) research programme.
Links and references to all relevant documents, presentations, code, and data
are provided where available.
\end{abstract}
%-----------------------------------------------------------------------------
%

{\color{white}Spacer}
\\[4cm]

% Add the license information (CC-BY-4).
\input{common/tools/licenseCCBY4}

\clearpage

% Add a table of contents.
\setcounter{tocdepth}{2}
\tableofcontents

\clearpage

% The introduction.
%%%%%%%%%%%%%%%%%%%%%%%%%%%%%%%%%%%%%%%%%%%%%%%%%%%%%%%%%%%%%%%%%%%%%%%%%%%%%%
\section{Introduction}
\label{sec:intro}
%%%%%%%%%%%%%%%%%%%%%%%%%%%%%%%%%%%%%%%%%%%%%%%%%%%%%%%%%%%%%%%%%%%%%%%%%%%%%%
CERN@school is a student-led research programme that aims to bring
\acs{CERN} into the classroom.
It is largely based around the
Timepix hybrid silicon pixel
detector~\cite{Timepix2007,TimepixErratum2008},
a device developed by the Medipix2 Collaboration\footnote{%
See \url{http://medipix.web.cern.ch}}
that enables the user to visualise and record measurements of
ionising radiation in real time.
As well as changing the way radiation can be explored in the classroom,
data from Timepix devices deployed around the UK, in space, and
at the \ac{LHC} can also be used by
CERN@school collaboration members to carry out authentic, novel scientific
research suitable for publication in peer-reviewed journals.
%

A brief history and description of the programme may be found in
the proceedings of the 37$\thsuper$ International Conference on
High Energy Physics (ICHEP 2014)~\cite{Whyntie2016b}.
%
Since that conference, held in July 2014, CERN@school has become a 
flagship programme of the
Institute for Research in Schools\footnote{%
See \url{http://researchinschools.org}} (IRIS),
a UK charity that aims to offer schools students and their teachers
the opportunity to participate in real science in many different fields.
%
This document aims to capture the progress and status of the
CERN@school programme in the period from June 2012
to December 2016.
This roughly corresponds to the start of the
\ac{STFC}
Science in Society Large Award (grant number ST/J000256/1)
and the end of the
\ac{STFC} \ac{PEF} for CERN@school
(grant number ST/N00101X/1).
%
To this end, the documents, presentations,
events, code repositories, and publications associated with
CERN@school research, educational, and promotional activities
have been indexed and are briefly described in the relevant
sections and subsections of this document.
%
Hyperlinks and \acp{DOI} are provided where
available.
%

%=============================================================================
\subsection{Overview of the index}
\label{sec:introoverview}
%=============================================================================
Section~\ref{sec:cas} indexes the
documents (Section~\ref{sec:casdocuments}),
presentations (Section~\ref{sec:caspresentations}),
events (Section~\ref{sec:casevents}),
and
code repositories (Section~\ref{sec:cascode})
related to the core CERN@school programme.
This includes research-based and educational material related to
the Timepix detector that has been produced as a direct result of the
programme, such as the
\acf{LUCID} %Langton Ultimate Cosmic ray Intensity Detector (LUCID)
satellite experiment,
\acf{RAY}, %Radiation Around You (RAY),
and
the \acl{UKSA}-backed TimPix Principia project.
%

Section~\ref{sec:moedal} provides the same information
(Sections~\ref{sec:moedaldocuments}--\ref{sec:moedalcode})
for items relating to the \ac{MoEDAL} experiment.
This material has been given its own section as
%CERN@school/IRIS joined MoEDAL, and
generally speaking the work reported linked to the \acs{STFC} \acs{PEF}
and was therefore more research-focussed.
%
The peer-reviewed publications produced as part of
CERN@school and/or \acs{MoEDAL} have been listed separately
in Section~\ref{sec:publications}.
%
References, acronyms, and acknowledgements may be found in
sections~\ref{sec:references},~\ref{sec:acronyms},
and~\ref{sec:ack} respectively.


\clearpage

% The CERN@school section

%%%%%%%%%%%%%%%%%%%%%%%%%%%%%%%%%%%%%%%%%%%%%%%%%%%%%%%%%%%%%%%%%%%%%%%%%%%%%%
\section{The CERN@school programme}
\label{sec:cas}
%%%%%%%%%%%%%%%%%%%%%%%%%%%%%%%%%%%%%%%%%%%%%%%%%%%%%%%%%%%%%%%%%%%%%%%%%%%%%%
As noted in the introduction, an overview of the CERN@school
programme may be found in the proceedings of the 2014
\ac{ICHEP}~\cite{Whyntie2016b}. A brief summary of the programme's
activities is provided below; the reader may then find further information
in the listings of
documents, %(Section~\ref{sec:casdocuments}),
presentations, %(Section~\ref{sec:caspresentations}),
events, %(Section~\ref{sec:casevents}),
and
code %(Section~\ref{sec:cascode})
(subsections~\ref{sec:casdocuments},
\ref{sec:caspresentations},
\ref{sec:casevents},
and~\ref{sec:cascode} respectively).
%
\begin{itemize}
\item \bullettext{In the beginning}:
CERN@school was conceived when students from the
Simon Langton Grammar School for Boys
visited the Medipix laboratories at \acs{CERN} in 2007.
%
After seeing the potential for using the
Timepix hybrid silicon pixel detectors~\cite{Timepix2007,TimepixErratum2008}
for detecting and visualising ionising radiation in the
classroom,
a pilot project was established with the support of
the \acf{SEPnet}
and
Kent County Council (KCC)
that placed ten detectors in schools around Kent, England.
%
\item \bullettext{LUCID}: 
Students from the Langton Star Centre,
the research facility attached to the school and
directed by B. Parker,
then submitted a proposal for the
\acf{LUCID}
experiment to a satellite competition for schools
run by the British National Space Centre\footnote{
  Now the \acf{UKSA}.
} and
\acf{SSTL}.
%
Consisting of five Timepix detectors arranged in an
open-faced cube, \ac{LUCID} was ultimately accepted as one of
scientific payloads aboard SSTL's \ac{LEO} technology demonstration
satellite, TechDemoSat-1.
%
In 2011, the UK \acf{STFC}
awarded \ac{LUCID} a Science in Society Large Award to support the
dissemination of data and the development of educational and
promotional resources for schools.
%
In the run-up to the (delayed) launch,
GEANT4~\cite{GEANT42003,GEANT42006}
simulations of LUCID were performed to estimate the performance
of the detector and provide training data for CERN@school students.
The results of this work were presented at the
\acs{CHEP} 2013~\cite{Whyntie2014} and \acs{iWoRID} 2014~\cite{Whyntie2015a}
conferences.
%
On Tuesday the 8th July 2014,
TechDemoSat-1 launched successfully aboard a Soyuz 2b
launch vehicle from Baikonur Cosmodrome, Kazakhstan. The satellite has
since been commissioned and has been delivering data for the CERN@school
programme since October 2014. LUCID project leader
C. Hewitt presented the first results at the 2015 \acf{NAM}
(see section~\ref{meeting:nam2015lucid}).
%
\item \bullettext{The Detector Network}:
With the launch of TechDemoSat-1 delayed
and additional funding secured,
the opportunity was taken 
to expand the original CERN@school detector network.
Twenty-five Timepix detectors,
supplied by Jablotron\footnote{%
  See \texttt{http://www.jablotron.com}},
were procured and distributed around the UK
via the \acf{IOP}
\acf{PTN} (see section~\ref{meeting:iop2013maydetdist}).
%
Many presentations and workshops were given to introduce
students, teachers, and outreach network coordinators
to the Timepix technology and data management systems.
%
In 2013, a Royal Commission for the Exhibition of 1851
Special Award was made to support the role of
Schools Research Champion for this national network of
detectors (and \ac{LUCID}).
%
Subsection~\ref{meeting:eclipse2015event} describes the
CERN@school Eclipse 2015 event, where many of the Timepix
detectors in schools made measurements before, during, and
after the total solar eclipse of Friday the 20th of March
2015~\cite{Whyntie2016a}.
%
Results from this experiment, and many others including
the \ac{RISE}, \ac{RAY}, and \ac{FAIR}, were presented at the
three CERN@school Research Symposia held at
the University of Surrey (2014, section~\ref{meeting:cassymposium2014}),
\acl{QMUL} (2015, section~\ref{meeting:cassymposium2015}),
and
the \acl{RAL} (2016, section~\ref{meeting:cassymposium2016}).
%
\item \bullettext{The demonstration experiments and user guides}:
as the Timepix detector is a research-grade piece of scientific
equipment, four classroom demonstration experiments were carried out,
written up and published in peer-reviewed journals -- {\em Physics
Education}~\cite{Whyntie2013} and
{\em Contemporary Physics}~\cite{Whyntie2015b}.
This was to demonstrate not only the physics but how to prepare
and publish one's scientific work and results with CERN@school.
User guides and instructions were also provided on the 
CERN@school website, which was at the time hosted by \acs{CERN}\footnote{%
Formerly at \href{http://cernatschool.web.cern.ch}{http://cernatschool.web.cern.ch}}.
The datasets from the experiments have also been made available
in the \href{http://fighare.com}{FigShare} online data repository
(see Section~\ref{sec:casdocuments}).
Following feedback from teachers and outreach officers,
a new set of guides and demonstration experiment worksheets were
prepared for \acs{IRIS} by LFT Consulting and the \acf{SSERC}.
CERN@school also moved its web presence to the \acs{IRIS} website,
where these documents are hosted (see Section~\ref{sec:casdocuments}).
%
\item \bullettext{Coding with CERN@school}:
The relatively simple data format used by the Timepix detector -- \code{ASCII}
text files containing the pixel $x$, $y$, and count values -- meant
that data processing and analysis with high-level programming languages like
Python could be done fairly easily.
CERN@school could therefore provide students (and teachers) with an
introduction to coding motivated by scientific research
(as opposed to making an app or other more trivial applications).
To this end, various training materials featuring CERN@school
data were prepared (see Section~\ref{sec:casdocuments}).
A CERN@school \href{http://github.com}{GitHub} educational
organisation was also created and CERN@school-related code was published there,
including the code used in the published CERN@school experiments.
See Section~\ref{sec:cascode} for further information and a list of
relevant code repositories.
%
\item \bullettext{GridPP and the CernVM}:
in October 2013, the Fellow joined the
GridPP Collaboration~\cite{GridPP2006,GridPP2009} as Dissemination
Officer.
%
GridPP, which represents the UK's contribution to the
\ac{WLCG}, had previously offered computing support for
CERN@school via the \texttt{cernatschool.org} \ac{VO}.
This appointment offered the opportunity to formalise this
partnership. \texttt{cernatschool.org} acted as a
technology demonstrator \ac{VO} for GridPP's
New User Engagement Programme.
By using these technologies
to power the work featured in CERN@school publications,
e.g.~\cite{Whyntie2015a}, they were tested and ready for
other new user communities
from medical physics, astrophysics and computational biology
to use in their activities and increasing GridPP's engagement
beyond the \ac{LHC} and particle physics.
Further details and case studies (of which CERN@school is one)
can be found in the document described in
Section~\ref{meeting:gridppcasestudies}.
CERN@school students were also involved in testing and providing
feedback on the GridPP UserGuide, the document described
in Section~\ref{meeting:gridppuserguide} which provides a
complete guide to getting on the Grid with GridPP.
Furthermore, the GridPP CernVM -- a \ac{VM} image provided
by \acs{CERN} configured for use with the GridPP New User
Engagement Programme -- can provide an out-of-the-box
working environment for coding with Python and
CERN@school. A document describing how to create a GridPP CernVM
for use with CERN@school is described in Section~\ref{meeting:cascernvmcheatsheet}.
CERN@school was also represented at many of the GridPP Collaboration
meetings in this period; see Section~\ref{sec:caspresentations}
for further details.

%
\item \bullettext{TimPix and Principia}:
following a successful partnership for a \acs{UKSA}-funded
\acs{LUCID} side project for primary schools,
``{\em Whatever the Space Weather}'',
L. F. Thomas (LFT Consulting) worked on behalf of
the Langton Star Centre and then the \ac{IRIS} to establish the
TimPix project.
Part of the \ac{UKSA} Principia Mission,
which coordinated the outreach activities associated with
UK \ac{ESA} astronaut Tim Peake's visit to the \ac{ISS},
TimPix gave participating students access to data from the
five Timepix detectors installed aboard the \ac{ISS}.
Students working on TimPix projects presented their work
at the Summer 2016 Medipix Collaboration meeting (Section~\ref{meeting:mpxcol2016cern})
and the Principia Schools Conference (Section~\ref{meeting:principia2016pompey}).
A CERN@school Timepix detector also featured in a demonstration
of real-time radiation visualisation in Dr Kevin Fong's
2015 Royal Institution Christmas Lectures (Section~\ref{meeting:rixmaslectures2015}).
%
\item \bullettext{The MoEDAL experiment at CERN}:
owing to the students' experience with the Timepix detector,
the Langton Star Centre joined the \ac{MoEDAL}
collaboration in 2013\footnote{%
See \href{https://cds.cern.ch/journal/CERNBulletin/2013/48/News Articles/1630194}{CERN Bulletin, Issue No. 48--49/2013}.},
and so CERN@school added data from a \acf{LHC} experiment
to its scientific portfolio\footnote{%
CERN@school@CERN, if you will.}.
However, as this became the focus of the Fellow's activities from
January 2015, the documents, presentations, events and code
repositories have been given a dedicated section
(Section~\ref{sec:moedal}).
\end{itemize}


%\clearpage


%=============================================================================
\subsection{Documents}
\label{sec:casdocuments}
%=============================================================================
Table~\ref{tab:casdocuments} lists the documents relating to
CERN@school research and activities,
which are then briefly described in the following subsections.
%
This includes things like internal reports, posters, experimental protocols,
published datasets, and user guides, but not
presentation slides (which are included with the
relevant presentation in Section~\ref{sec:caspresentations})
or code repositories (which are described in Section~\ref{sec:cascode}).
%
Peer-reviewed publications from the CERN@school programme have been
listed in Section \ref{sec:publications} and are marked with an
asterisk in Table~\ref{tab:papers1}.

\begin{landscape}

\input{cas/documents/documentstable}

\end{landscape}

\input{cas/documents/documents}

\clearpage

%=============================================================================
\subsection{Presentations}
\label{sec:caspresentations}
%=============================================================================
Table~\ref{tab:caspresentations} provides a list of selected presentations
relating to the CERN@school programme from June 2012 to December 2016,
i.e. the period covering the STFC Large Award, the Special Award from
the Royal Commission for the Exhibition of 1851, and the
CERN@school STFC Public Engagement Fellowship.
These are then detailed in the following subsections,
except for standard school visits (i.e. talks in schools where an overview of
the programme and basic detector and data training was provided).
Note that MoEDAL-related presentations are featured in
Table~\ref{tab:moedalmeetings} (Section~\ref{sec:moedalpres}).

\input{cas/presentations/presentationstable}

\clearpage

\input{cas/presentations/presentations}

\clearpage

%=============================================================================
\subsection{Events}
\label{sec:casevents}
%=============================================================================
Table~\ref{tab:casevents} lists events that took place between
June 2012 and January 2017 that in some way involved CERN@school.
Each event is then briefly described in the following subsections.
Please note that,
generally speaking,
if there was a presentation given at an event
it will be listed in Section~\ref{sec:caspresentations}
and may not be listed here.

\input{cas/events/eventstable}

\clearpage

\input{cas/events/events}

\clearpage

%=============================================================================
\subsection{Code repositories}
\label{sec:cascode}
%=============================================================================
In order to ensure students could freely participate in
CERN@school-related scientific studies,
all software and analysis code required to analyse data from the Timepix
detectors was made publicly available via the online code repository
\href{http://github.com}{GitHub}.
When CERN@school existed as a separate entity,
CERN@school used the CERN@school GitHub Organization:

\href{http://github.com/CERNatschool}{http://github.com/CERNatschool}

as a place for hosting software repositories that were undergoing active development.
GitHub kindly provided an Educational Account for CERN@school
which meant that a limited number of private repositories could be kept too.
%
However, with the creation of the \acf{IRIS}, a new GitHub Organization
was created for software development across all \ac{IRIS} projects:

\href{http://github.com/InstituteForResearchInSchools}{href://github.com/InstituteForResearchInSchools}

and the CERN@school Organization has since been decommissioned.
As a number of the repositories featured code used
in the CERN@school documents and publications,
the repository will remain on the GitHub website but will no longer
be active.
Except where otherwise noted,
the software is issued under an MIT software license
and may be used by \acs{IRIS} students (or indeed anyone else)
as required.
Furthermore, stable releases of these repositories
have been published on the \acs{CERN}-backed
\href{http://zenodo.org}{Zenodo} repository
which assigns them a permanent \ac{DOI}.
The relevant repositories are listed in
Table~\ref{tab:cascode} and detailed in the
following subsections.

\input{cas/code/codetable}

\clearpage

\input{cas/code/code}

\clearpage

\clearpage

% The MoEDAL section.

%%%%%%%%%%%%%%%%%%%%%%%%%%%%%%%%%%%%%%%%%%%%%%%%%%%%%%%%%%%%%%%%%%%%%%%%%%%%%%
\section{MoEDAL}
\label{sec:moedal}
%%%%%%%%%%%%%%%%%%%%%%%%%%%%%%%%%%%%%%%%%%%%%%%%%%%%%%%%%%%%%%%%%%%%%%%%%%%%%%
The \acl{MoEDAL}~\cite{MoEDAL2009}\
is the seventh major experiment at \acs{CERN}'s \acf{LHC}.
It is housed underground at Interaction Point (IP8), sharing
the experimental cavern with the \acs{LHCb} experiment.
%
\acs{MoEDAL} is designed to probe the \acs{LHC}'s particle collisions
for signs of Paul Dirac's hypothesised magnetic monopole,
as well as other highly ionising signs of new physics that
the other \acs{LHC} experiments cannot easily look for.
%
A full introduction to the \acs{MoEDAL} experiment may be found
in document CAS-PUB-MDL-000002 (see Section~\ref{meeting:moedalcasguide})
and an accompanying bibliography in CAS-PUB-MDL-000001
(see Section~\ref{meeting:moedalbib}); the reader is advised to
consult these for more information about \ac{MoEDAL} itself.
%

The \acf{IRIS} is a full member of 
the \acs{MoEDAL} Collaboration, with the Langton Star Centre having
joined in 2013 owing to their experience with the Timepix detectors.
Students and teachers associated with \acs{IRIS}
have access data from the \acs{MoEDAL} experiment and are able to
contribute to the \acs{MoEDAL} Collaboration's research programme.
%
With the expansion of the CERN@school team under \acs{IRIS},
it was decided that from January 2016
the Fellow should focus on research activities
with the \acs{MoEDAL} Collaboration. These included:

\begin{itemize}
\item \bullettext{Collaboration activities}:
as with any \acs{LHC} experiment, there were many collaboration activities
that the Fellow was required to perform as part of the role.
These included the biannual Collaboration Meetings
and various trips to \acs{CERN}.
As a member of the \ac{SaAG},
there were also the \ac{SaAG} meetings every other week.
All of these meetings, with brief summaries,
are listed in the Presentations and/or Events sections.
Furthermore, there were various promotional activities
undertaken on behalf of the Collaboration. These included the
Monopole Quest! exhibit at the Royal Society Summer Science Exhibition
(see Section~\ref{meeting:rssse2015}),
various schools talks,
and an article published via The Conversation which was read by
over half a million people (see Section~\ref{meeting:moedalconv}).
%
\item \bullettext{The Timepix detector network}:
funding for two Timepix detectors for deployment in the \acs{LHCb}
cavern was secured from \acs{STFC}. This deployment has been managed by
the University of Alberta on behalf of CERN@school (and so \acs{IRIS}).
As full collaboration members, CERN@school also has access to 
the Timepix detector network managed by the \acf{CTU}, Prague.
Data from this network, which was operational during Runs I and II of
\acs{LHC} operations, has been used as the basis of studies by
\acs{IRIS} and \acs{MoEDAL} \acs{CERN} summer school students which
forms the basis of the \acs{MoEDAL} internal note,
``{\em On the flora and fauna of highly ionising background radiation
at Interaction Point 8}''.
Further details may be found in Section~\ref{meeting:moedaltimepixnote}.
%
\item \bullettext{The Nuclear Track Detectors and Monopole Quest!}:
as with all research programmes, new developments and opportunities
led to a slight shift in direction for the Fellow's activities.
While CERN@school's expertise originally focussed on the Timepix detector,
it was suggested that the thousands of visitors to the \ac{RSSSE} 
could be engaged with \acs{MoEDAL}'s research by helping to analyse
scanned images from the \acp{NTD}.
With this in mind,
a Zooniverse Citizen Science project was designed and created
by \acs{IRIS} students using the
Panoptes project builder. This enabled volunteers (at the \ac{RSSSE}
or elsewhere)
to identify and measure etch pits and other features in scans
from \ac{NTD} plastic exposed to the \acs{LHC} radiation environment
and heavy ion test beams.
The results from this formed the basis of a pilot study
reported in a \acs{MoEDAL} internal note
(see Section~\ref{meeting:moedalntdnote}).
Furthermore, a guide to MoEDAL (CAS-PUB-MDL-000002)
and accompanying educator presentation template (CAS-PUB-MDL-000004)
were prepared to facilitate the use of Monopole Quest! in
a classroom context. Details of these can be found in
Sections~\ref{meeting:moedalcasguide} and~\ref{meeting:moedalschtemplate}
respectively.
%
\item \bullettext{MoEDAL software and the Grid}:
towards the end of 2016, the Fellow also established the
\acs{MoEDAL} \acf{VO} on the \acf{WLCG} which enabled both
\acs{MoEDAL} and CERN@school researchers to make use of
GridPP's computing resources.
Following this, and a successful run of simulations used in the
analysis of 13 TeV data from the
\acf{MMT} subdetector system,
the Fellow was appointed Software Coordinator for the
\acs{MoEDAL} Collaboration.
As part of the duties associated with this role,
the MoEDAL Grid User Guide (CAS-PUB-MDL-000006)
and \code{moedal-run-simulations} repository
were created; these are detailed in Sections~\ref{meeting:moedalgriduserguide}
and~\ref{sec:moedalrunsimulations} respectively.
MoEDAL also featured as a case study for
GridPP's New User Engagement Programme.
MoEDAL (and \acs{LHCb}) software is available to
\acs{IRIS} students via \acs{CVMFS} on a GridPP CernVM.
\end{itemize}

As a member of the \acs{MoEDAL} Collaboration,
the Fellow was also included on the \acs{MoEDAL} author list.
Associated publications are listed in Section~\ref{sec:publications}.
These include the first physics results from the \acs{MoEDAL}
experiment as reported in the \acf{JHEP}~\cite{MoEDAL2016a}
and \acf{PRL}~\cite{MoEDAL2017a}.
The reader may also be interested in
the CERN@school bibliography (CAS-PUB-MDL-000001),
which contains a fairly comprehensive set of references
relating to MoEDAL and many other searches for magnetic monopoles.
Further details may be found in Section~\ref{meeting:moedalbib}.

\clearpage

%=============================================================================
\subsection{Documents}
\label{sec:moedaldocuments}
%=============================================================================
Table~\ref{tab:moedaldocuments} lists the documents relating to
\acs{MoEDAL} research and activities supported by \acs{IRIS}.
A brief description of each is provided 
in the following subsections.

\input{moedal/documents/documentstable}

\input{moedal/documents/documents}

\clearpage

%=============================================================================
\subsection{Presentations}
\label{sec:moedalpres}
%=============================================================================
Table~\ref{tab:moedalmeetings} lists the presentations given as part
of the Fellow's work with the \acs{MoEDAL} Collaboration.
Each one is then briefly described in the following
subsections.
%
Please note that the \acs{MoEDAL}
\acf{SaAG} % Software and Analysis group meetings
and \acs{MoEDAL} Collaboration meetings are only available to
\acs{MoEDAL} Collaboration members with the appropriate
credentials.
%
If access is required, please contact the
\acl{IRIS}\footnote{%
\url{http://researchinschools.org}}.

\input{moedal/presentations/presentationstable}

\clearpage

\input{moedal/presentations/presentations}

\clearpage

%=============================================================================
\subsection{Events}
\label{sec:moedalevents}
%=============================================================================
Table~\ref{tab:moedalevents} lists MoEDAL-related events that
have taken place since CERN@school first became involved with the
experiment.
The events themselves are then briefly described in the following subsections.
Note that for MoEDAL, this includes school visits, student workshops,
and other outreach and engagement-focussed events.
%
It should also be noted that the first MoEDAL Collaboration meeting
(also held at CERN) was not officially a collaboration meeting,
being more of a monopole physics workshop.  Regardless, it occured
before CERN@school was involved with the experiment so is not included
in the list of MoEDAL/CERN@school events.

\input{moedal/events/eventstable}

\clearpage

\input{moedal/events/events}

\clearpage

%=============================================================================
\subsection{Code repositories}
\label{sec:moedalcode}
%=============================================================================
The shared code repositories used for work on the \acs{MoEDAL} experiment
are listed in Table~\ref{tab:moedalcode}.
The code required to generate various public guides and documents 
has been archived on the CERN@school GitHub repository:

\href{http://github.com/CERNatschool}{http://github.com/CERNatschool}

Active analysis code and MoEDAL data cannot reside in the public domain
(at least until the associated results have been published by the
\acs{MoEDAL} collaboration). The corresponding repositories are therefore
hosted on \acs{CERN}'s GitLab system available to those with \acs{CERN}
computing accounts. \acs{MoEDAL}'s repositories are hosted here:

\href{https://gitlab.cern.ch/moedal/}{https://gitlab.cern.ch/moedal/}

Students and teachers who have got as far as requiring access to the
non-public MoEDAL data should contact \acs{IRIS} about obtaining a
\acs{CERN} computing account and joining the \acs{MoEDAL}
\acl{SaAG}.

\input{moedal/code/codetable}

\input{moedal/code/code}

\clearpage

\clearpage


%%%%%%%%%%%%%%%%%%%%%%%%%%%%%%%%%%%%%%%%%%%%%%%%%%%%%%%%%%%%%%%%%%%%%%%%%%%%%%
\section{Publications}
\label{sec:publications}
%%%%%%%%%%%%%%%%%%%%%%%%%%%%%%%%%%%%%%%%%%%%%%%%%%%%%%%%%%%%%%%%%%%%%%%%%%%%%%
The peer-reviewed publications associated with and/or referenced by
the CERN@school programme are listed with full bibliographic citations in the
\hyperref[sec:references]{References section}.
%
For convenience,
the \code{BibTeX} citation code used in this document's \code{BibTeX} file,
DOI or URL, brief notes, year of publication and citation number
are listed in Table~\ref{tab:papers1}.
%
Papers that were directly supported by either
STFC grant number ST/J000256/1
(Science in Society Large Award, plus extensions)
or
STFC grant number ST/N00101X/1
(Public Engagement Fellowship)
are marked with an asterisk.

\begin{landscape}
%
% The table produced by make_paper_summary_tables.py
\input{autotable}
%
\end{landscape}

\newpage

%
%%%%%%%%%%%%%%%%%%%%%%%%%%%%%%%%%%%%%%%%%%%%%%%%%%%%%%%%%%%%%%%%%%%%%%%%%%%%%%%
% Bibliography
%%%%%%%%%%%%%%%%%%%%%%%%%%%%%%%%%%%%%%%%%%%%%%%%%%%%%%%%%%%%%%%%%%%%%%%%%%%%%%%
%
\section{References}
\label{sec:references}
\bibliographystyle{unsrt.bst}
\bibliography{cernatschool}
%
%------------------------------------------------------------------------------

\clearpage

%%%%%%%%%%%%%%%%%%%%%%%%%%%%%%%%%%%%%%%%%%%%%%%%%%%%%%%%%%%%%%%%%%%%%%%%%%%%%%
\section{Acronyms}
\label{sec:acronyms}
%%%%%%%%%%%%%%%%%%%%%%%%%%%%%%%%%%%%%%%%%%%%%%%%%%%%%%%%%%%%%%%%%%%%%%%%%%%%%%

\begin{acronym}
\acro{APPE}{Advisory Panel for Public Engagement}
\acro{CERN}{Organisation Europ\'een pour la Recherche Nucl\'aire}
\acro{CHEP}{\acroextra{The International Conference on }Computing in High Energy and Nuclear Physics}
\acro{CTU}{Czech Technical University\acroextra{ in Prague}}
\acro{CVMFS}{CernVM File Service}
\acro{DOI}{Digital Object Identifier}
\acro{ESA}{European Space Agency}
\acro{FAIR}{Fractional Attenuation of Ionising Radiation\acroextra{ (CERN@school experiment)}}
\acro{ICHEP}{International Conference of High Energy Physics}
\acro{IFIC}{Instituto de Fisica Corpuscular}
\acro{IOP}{Institute of Physics}
\acro{IRIS}{\acroextra{The }Institute for Research in Schools}
\acro{ISS}{International Space Station}
\acro{iWoRID}{International Workshop on Radiation Imaging Detectors}
\acro{JHEP}{Journal of High Energy Physics}
\acro{KCL}{King's College London}
\acro{LEO}{Low Earth Orbit}
\acro{LHC}{Large Hadron Collider}
\acro{LHCb}{\acroextra{The }Large Hadron Collider Beauty\acroextra{ experiment}}
\acro{LUCID}{Langton Ultimate Cosmic ray Intensity Detector}
\acro{MoEDAL}{Monopole and Exotics Detector at the LHC}
\acro{MMT}{Magnetic Monopole Trapper\acroextra{ (MoEDAL subdetector)}}
\acro{NAM}{National Astronomy Meeting}
\acro{NASA}{National Aeronautics and Space Administration}
\acro{NTD}{Nuclear Track Detector\acroextra{ (MoEDAL subdetector)}}
\acro{PEF}{Public Engagement Fellowship}
\acro{PNC}{Physics Network Coordinator\acroextra{ (c.f. PTN)}}
\acro{PPRC}{\acroextra{The QMUL }Particle Physics Research Centre}
\acro{PRL}{Physical Review Letters}
\acro{PTN}{\acroextra{The IOP }Physics Teacher Network}
\acro{QMUL}{Queen Mary University of London}
\acro{RAL}{\acroextra{The }Rutherford Appleton Laboratory}
\acro{RAY}{Radiation Around You\acroextra{ (CERN@school experiment)}}
\acro{RISE}{Radiation In Schools Experiment\acroextra{ (CERN@school experiment)}}
\acro{RSSSE}{Royal Society Summer Science Exhibition}
\acro{SaAG}{\acroextra{The MoEDAL }Software and Analysis Group}
\acro{SEPnet}{South East Physics network}
\acro{SLC}{Science Learning Centre}
\acro{SSERC}{Scottish Schools Education Research Centre}
\acro{SSTL}{Surrey Satellite Technology Limited}
\acro{STFC}{Science and Technology Facilities Council}
\acro{UKSA}{UK Space Agency}
\acro{VM}{Virtual Machine}
\acro{VO}{Virtual Organisation}
\acro{WLCG}{Worldwide LHC Computing Grid}
\end{acronym}


%\clearpage

%%%%%%%%%%%%%%%%%%%%%%%%%%%%%%%%%%%%%%%%%%%%%%%%%%%%%%%%%%%%%%%%%%%%%%%%%%%%%%%
\section{Acknowledgements}
\label{sec:ack}
%%%%%%%%%%%%%%%%%%%%%%%%%%%%%%%%%%%%%%%%%%%%%%%%%%%%%%%%%%%%%%%%%%%%%%%%%%%%%%%
CERN@school has grown from a single school's visit to \acs{CERN}
into a nationwide flagship programme for the \acf{IRIS},
via a regional pilot scheme and national expansion supported by
\ac{SEPnet} and \ac{STFC} respectively.
This would not have been possible without the support of a great number of
people and organisations who have seen the potential of the programme
and what its philosophy offers participating students and teachers.
While the following list is not exhaustive, it is hopefully indicative of
the tremendous amount of help that has been offered and provided
by scientists, experts, and educational practioners from around the world.

Firstly, there would be no CERN@school without the
Medipix Collaboration. The Timepix detector is a phenomenal piece of
research-grade equipment that has made, and will continue to make,
a huge impact in many fields
that require real-time visualisation of ionising radiation.
Science education is undoubtedly one of these fields.
M. Campbell, spokesperson for the Medipix Collaboration,
has provided a huge amount of support from day one
and has been, quite literally, instrumental in its success.
We also thank J. A. Alozy for support with detector repair
and maintanence, J. Id\'{a}ragga, M. Benoit and S. Arfaoui
for the GEANT4-based AllPix simulation suite used
in~\cite{Whyntie2014,Whyntie2015a,Whyntie2015b},
and the whole of the Medipix Collaboration for their
hospitality and engagement at various collaboration meetings.
%
CERN, likewise, has been very supportive of the programme.
We are grateful for the continued permission to use the
\acs{CERN} name and logo for CERN@school activities and research,
and would like to thank J. Gillies, R. Vanden Broeck, and K. Kahle
for their help engaging with CERN and beyond.
The author would also like to thank the Users Office,
IT Department, Hostel service, and the staff in Restaurant 1
for providing a consistently excellent working environment\footnote{%
And exceptional croissants fourr\'{e} cr\`{e}me p\^{a}tissi\`{e}re.}.
There truly is nowhere like it in the world.
%

The support and guidance of the UK \acf{STFC} has been essential to
the sustained growth of CERN@school, both in terms of the team and
equipment. The \acs{STFC} Public Engagement group's
Science in Society programme and Public Engament Strategy
have supported CERN@school through the very successful Large Award and
\acf{PEF} schemes, as well as additional
Capital Resource allocations.
Furthermore, the guidance and oversight provided by the
\acs{STFC} Steering Group for CERN@school ensured that the programme
ran as it should:
we would like to thank
R.~Clegg,
E.~Cunningham,
D.~A.~Gillespie,
N.~Hollingworth,
A.~Thompson,
and
C.~Woolford
from the \acs{STFC} Public Engagement team,
and
J.~P.~Eastwood,
C.~Johnson, % CHECK
C.~P.~R.~Thorley,
and the various educators who joined the Steering Group for this.
A special note of thanks should also go to
J.~Womersley, % CHECK 
former Chief Executive of \acs{STFC}, % CHECK
for supporting CERN@school so enthusiastically during his time in the role.
%
Likewise, the contribution of the
Royal Commission for the Exhibition of 1851 via a Special Award
ensured funding for the Fellow continued after the Large Award ended and
the \ac{PEF} began; thanks are due to N.~Williams and
his team with respect to this.

The Langton Star Centre,
and the Simon Langton Grammar School for Boys\footnote{%
Girls are admitted in the sixth form.},
provided an ideal environment for CERN@school to grow and flourish.
%
The support, freedom, and trust students are given to pursue their own
lines of investigation their have been paramount to the success of the
student-led projects that have started there.
However, this would not be possible without a fantastic team of
staff working tirelessly behind the scenes to make sure the
research and educational needs of the students are met.
The author is grateful to the following for their help:
R.~Champion, T.~Connolly (educational support);
S.~Begg, J.~Bradley, S.~Greenwood,
N.~Hatton, H.~Mudhar, E.~Mudhar-Fung,
J.~Poyner, and D.~Wash (adminstrative support);
I.~Begg, C.~Boucher, T.~Forester, and M.~Kemp (technical support);
and
M.~N.~F.~Baxter and K.~Moffat (Senior Management Team).
%
A special note of thanks must go to G. Poole who led the
Science Department during the author's time at the school,
and provided invaluable guidance and the secure educational
foundations upon which the Langton Star Centre (and so CERN@school)
could be built.

The CERN@school network could not have expanded as it did without
two key organisations: the \acf{SEPnet} and the UK \acf{IOP}.
The former successfully piloted a regional programme via
their university-based outreach officers, while the latter
expanded the network across the UK and Ireland through their
\acf{PTN}.
We thank
C.~Harvey
and
J.~West
of \acs{SEPnet}
and
C.~Shepherd
and
G.~Williams
of the \acs{IOP}
for their continued support and expertise.
V.~Benson (\acs{SEPnet})
and
V.~Fox (\acs{IOP})
also provided a number of exceptional summer students via their
respective internship programmes.
The contributions of A.~Coupe, J.~Cook, and N.~Shearer
have been noted elsewhere in this document;
they were a pleasure to supervise and the author wishes them
all the best in their future careers.
%
There are too many outreach officers and \acfp{PNC}
to mention by name, but likewise their contributions have been
greatly appreciated.
%
The author also wishes to acknowledge the support of
A.~Newsam (Liverpool John Moores University)
provided through the \acs{IOP} mentorship scheme.

Many of CERN@school's sub-projects have benefitted
from expert advice and guidance. These include:
\bullettext{LUCID} -- D.~Cooke, E.~Brownbill, S.~Ghanoun,
V.~O'Donovan, and S.~Wokes (\acs{SSTL}),
D.~J.~Garton (Satellite Applications Catapult),
J.~Jak\r{u}bek (\acs{CTU} Prague, detector calibration),
and J.~P.~Eastwood (Imperial College London);
the \bullettext{detector network} -- V.~Stanislav (Jablotron),
J.~R.~Wilson, C.~P.~R.~Thorley, and B.~Frost (\acs{QMUL}),
H.~Pollard (The Ogden Trust),
I.~Holmes (Kettering Buccleuch Academy),
S.~Palmer, L.~Mowberry, V.~Stowell, and J.~Lewis (\acs{RAL} Public Engagement),
L.~F.~Thomas (LFT Consulting) and G.~Steele (\acs{SSERC});
\bullettext{TimPix} -- L.~S.~Pinsky (\acs{NASA}/Uni. of Houston),
N.~N.~Stoffle (\acs{NASA}), Son Hoang (Uni. Houston),
L.~F.~Thomas (LFT Consulting -- again!),
L.~Jackson, J.~Curtis, and H.~Garrett (\acs{UKSA}),
and T.~Peake (\acs{ESA}).

Involvement with the \acs{MoEDAL} experiment started as a sub-project,
but as the \acs{STFC} \acs{PEF} (and CERN@school as a whole) evolved
so did the opportunities for extending the scope of the research activities.
This is largely down to the support provided by the \acs{MoEDAL}
Collaboration as led by J.~Pinfold.
Many people have provided help in the following areas:
R.~Soluk (Uni. Alberta) -- Timepix detector installation;
D.~Felea (ISS, Bucharest) -- simulations and detector geometry;
S.~Posp\'{i}\v{s}il, B.~Bergmann, and P.~Bene\v{s} (\acs{CTU} Prague) -- the
\acs{CTU} Timepix detector network;
and
V.~Wedlake (\acs{CERN}) -- adminisitrative support.
Special thanks should also go to
P.~Mermod (Uni. Geneva) for successfully chairing the
\acl{SaAG} and leading the efforts that
led to \acs{MoEDAL}'s first physics publications from the
\acs{MMT} subdetector results~\cite{MoEDAL2016a,MoEDAL2017a},
and to the various CERN Summer students who contributed
to \acs{IRIS} \acs{MoEDAL} projects.

On the computing front, the author would like to thank
the CernVM group (\acs{CERN}) for developing the CernVM
and CernVM-FS services that have made access to Grid resources
much, much easier for schools and other new user communities.
%
From GridPP, the following people have provided invaluable support
for CERN@school:
C.~Condurache (RAL) -- CernVM-FS;
%has supported GridPP's CernVM-FS
%capabilities at the Rutherford Appleton Laboratory (RAL);
%
D.~Bauer, S.~Fayer (Imperial College London),
and J.~Martyniak (now at RAL) -- GridPP DIRAC;
%developed and supporting GridPP DIRAC;
%
the Ganga development team\footnote{%
See \href{https://github.com/ganga-devs/}{https://github.com/ganga-devs/}};
%for their work on and support for the Ganga user interface;
S.~Skipsey (Uni. Glasgow) and J.~Jensen (RAL) -- data storage;
%from the
%GridPP Storage group for VO storage-related support;
%
A.~McNab (Uni. Manchester) -- web support;
N.~O'Neill (previously at QMUL) -- outreach;
%with technical support with
%the GridPP website and online user guide\footnote{%
%See \href{https://www.gridpp.ac.uk/userguide/}{https://www.gridpp.ac.uk/userguide/}};
%
D.~Taylor, A.~Owen, C.~Timis, T.~Froy, and C.~J.~Walker (QMUL) -- QMUL Tier-2 cluster support;
%support with development work on the Queen Mary University of London cluster;
%
J.~Coles (Uni. Cambridge) -- new user community support;
and
%for leading GridPP Operations and
%coordinating new user community activities;
%
P.~Gronbech (Uni. Oxford) -- project management.
%for support with Project Management;
%
%
Special thanks also go to
D.~Britton (Uni. Glasgow) for leading the GridPP Collaboration
and supporting the vision of integrating CERN@school via
the New User Engagement Programme
as it developed,
and
%
S.~L.~Lloyd (QMUL) for chairing the GridPP Collaboration Board,
providing numerous testing and support mechanisms for Grid
activities, initiating CERN@school's involvement with GridPP,
and providing much support and guidance as the two
organisations worked together.

Last, and by no means least, the author would like to thank
B.~Parker for leading CERN@school, the Langton Star Centre,
\acs{IRIS}, and the many other projects that have inspired
huge numbers of students in many areas of science and beyond.
It has been an incredible experience to work with someone the
author will always describe as the
``fifth fundamental force of Nature'',
and the impact she has made on the countless students,
educators, and indeed scientists through her work is
%(and perhaps always will be)
immeasurable.

{\em The work described here was supported by the
UK \acf{STFC} via grant numbers ST/J000256/1 and ST/N00101X/1,
and a Special Award from the Royal Commission for the Exhibition of 1851.
%
The CERN@school Collaboration would also like to acknowledge the
support provided by
the GridPP Collaboration~\cite{GridPP2006,GridPP2009}
in terms of both computing resources and technical guidance
from collaboration members.}


%\clearpage

%%%%%%%%%%%%%%%%%%%%%%%%%%%%%%%%%%%%%%%%%%%%%%%%%%%%%%%%%%%%%%%%%%%%%%%%%%%%%%%
\section*{Version History}
%%%%%%%%%%%%%%%%%%%%%%%%%%%%%%%%%%%%%%%%%%%%%%%%%%%%%%%%%%%%%%%%%%%%%%%%%%%%%%%
%______________________________________________________________________________
\begin{table}[h]
\caption{\label{tab:version}Version history.}
\lineup
\begin{indented}
\item[]\begin{tabular}{@{}cllc}
\br
\centre{1}{$\quad$Version    $\quad$} & 
\centre{1}{$\quad$Description$\quad$} &
\centre{1}{$\quad$DOI        $\quad$} &
\centre{1}{$\quad$Author     $\quad$} \\
\mr
1.0 & Initial version. & \href{http://doi.org/10.5281/zenodo.227090}{10.5281/zenodo.227090}& TW \\
\br
\end{tabular}
\end{indented}
\end{table}
%______________________________________________________________________________

\end{document}
